\subsection{线性组合引理}
我们将以证明每个矩阵都行等价于一个
而且仅有一个简化阶梯形矩阵来结束本章。
这里的思想将在下一章中再次出现，
并得到进一步的发展。
% Here is an informal argument that the reduced 
% echelon form version of a matrix is unique.
% Consider again the example that started this section of a matrix that
% reduces to three different echelon form matrices.
% The first matrix of the three is the natural echelon form version.
% The second matrix is the same as 
% the first except that a row has been halved.
% The third matrix, too, is just a cosmetic variant of the first. 
% The definition of reduced echelon form outlaws this kind of fooling around.
% In reduced echelon form,
% halving a row is not possible because that would
% change the row's leading entry away from one, and
% neither is combining rows possible, because then a leading entry would no
% longer be alone in its column.

% This informal justification is not a proof;
% the argument shows that no two different reduced echelon form matrices
% are related by a single row operation step, but the argument does not
% rule out the possibility that two different reduced echelon form
% matrices could be related by multiple steps.
% Before we go to the proof, we finish this subsection by 
% rephrasing our work in a terminology that will be enlightening.
关键的观察涉及
行变换如何把一个矩阵变成另一个矩阵：
新的行是
原来的行的线性组合。

\begin{example}
考察这个高斯–若尔当消元。
\begin{align*}
  \begin{amat}{2}
    2  &1  &0  \\
    1  &3  &5
  \end{amat}
  &\grstep{-(1/2)\rho_1+\rho_2}
  \begin{amat}{2}
    2  &1    &0  \\
    0  &5/2  &5
  \end{amat}                           \\
  &\grstep[(2/5)\rho_2]{(1/2)\rho_1}
  \begin{amat}{2}
    1  &1/2  &0  \\
    0  &1    &2
  \end{amat}                        \\
  &\grstep{-(1/2)\rho_2+\rho_1}\;
  \begin{amat}{2}
    1  &0    &-1  \\
    0  &1    &2
  \end{amat}                       
\end{align*}
将这些矩阵记为 $A\rightarrow D\rightarrow G\rightarrow B$，
并把 $A$ 的行写作 $\alpha_1$ 与 $\alpha_2$，等等，我们就有下式。
\begin{align*}  % multline looks worse
  \left(\begin{array}{l}
    \alpha_1 \\
    \alpha_2
  \end{array}\right)
  &\grstep{-(1/2)\rho_1+\rho_2}
  \left(\begin{array}{l}
    \delta_1=\alpha_1 \\
    \delta_2=-(1/2)\alpha_1+\alpha_2
  \end{array}\right)                          \\
  &\grstep[(2/5)\rho_2]{(1/2)\rho_1}
  \left(\begin{array}{l}
    \gamma_1=(1/2)\alpha_1 \\
    \gamma_2=-(1/5)\alpha_1+(2/5)\alpha_2
  \end{array}\right)                          \\
  &\grstep{-(1/2)\rho_2+\rho_1}
  \left(\begin{array}{l}
    \beta_1=(3/5)\alpha_1-(1/5)\alpha_2  \\
    \beta_2=-(1/5)\alpha_1+(2/5)\alpha_2
  \end{array}\right)                         
\end{align*} 
\end{example}

\begin{example}
高斯操作对行作线性组合这一事实，
在含有行对换时同样成立。
对于这里的 \( A \)、\( D \)、\( G \) 与 \( B \)
\begin{equation*}
    \begin{mat}[r]
       0  &2  \\
       1  &1
     \end{mat}
    \grstep{\rho_1\leftrightarrow\rho_2}
    \begin{mat}[r]
       1  &1  \\
       0  &2
     \end{mat}                  
    \grstep{(1/2)\rho_2}
    \begin{mat}[r]
       1  &1  \\
       0  &1
     \end{mat}                   
    \grstep{-\rho_2+\rho_1}
    \begin{mat}[r]
       1  &0  \\
       0  &1
     \end{mat}
\end{equation*}
我们得到这些线性关系。
\begin{multline*}
  \left(\begin{array}{l}
    \vec{\alpha}_1 \\
    \vec{\alpha}_2
  \end{array}\right)
  \,\grstep{\rho_1\leftrightarrow\rho_2}\,
  \left(\begin{array}{l}
     \vec{\delta}_1 =\vec{\alpha}_2   \\
     \vec{\delta}_2 =\vec{\alpha}_1
  \end{array}\right)                                                 
  \,\grstep{(1/2)\rho_2}\,
  \left(\begin{array}{l}
     \vec{\gamma}_1 =\vec{\alpha}_2  \\
     \vec{\gamma}_2 =(1/2)\vec{\alpha}_1
  \end{array}\right)                                               \\
  \,\grstep{-\rho_2+\rho_1}\,
  \left(\begin{array}{l}
     \vec{\beta}_1 =(-1/2)\vec{\alpha}_1+1\cdot\vec{\alpha}_2   \\
     \vec{\beta}_2 =(1/2)\vec{\alpha}_1
  \end{array}\right)
\end{multline*}
\end{example}

总之，
高斯消元法系统地找出行的一个合适的
线性组合序列。

% \begin{definition}
% A \definend{linear combination}\index{linear combination} of 
% \( x_1,\ldots,x_m \)
% is an expression of the form $c_1x_1+c_2x_2+\,\cdots\,+c_mx_m$
% where the \( c \)'s are scalars.
%\end{definition}

% \noindent (We have already used the phrase 
% `linear combination' in this book.
% The meaning is unchanged, but the next result's statement makes
% a more formal definition in order.)

\begin{lemma}[线性组合引理] \label{lm:LinearCombinationLemma} \index{Linear Combination Lemma}
%<*lm:LinearCombinationLemma>
线性组合的线性组合仍是线性组合。
%</lm:LinearCombinationLemma>
\end{lemma}

\begin{proof}
%<*pf:LinearCombinationLemma>
给定由 $x$ 们的线性组合组成的集合，
从 $c_{1,1}x_1+\dots+c_{1,n}x_n$ 到 $c_{m,1}x_1+\dots+c_{m,n}x_n$，
考虑这些组合的一个组合
\begin{equation*}
  d_1(c_{1,1}x_1+\dots+c_{1,n}x_n)\,+\dots+\,d_m(c_{m,1}x_1+\dots+c_{m,n}x_n)
\end{equation*}
其中 $d$ 们与 $c$ 们一样都是标量。
将这些 $d$ 分配开来并重新组合，得到
\begin{equation*}
  %&=d_1c_{1,1}x_1+\dots+d_1c_{1,n}x_n\,
  % +d_2c_{2,1}x_1+\dots+\,
  % d_mc_{1,1}x_1+\dots+d_mc_{1,n}x_n         \\
  =(d_1c_{1,1}+\dots+d_mc_{m,1})x_1\,+\dots+\,(d_1c_{1,n}+\dots+d_mc_{m,n})x_n
\end{equation*}
它仍然是 $x$ 们的一个线性组合。
%</pf:LinearCombinationLemma>
\end{proof}


\begin{corollary} \label{cor:RowsOfEqMatsLinCombos}
%<*co:RowsOfEqMatsLinCombos>
当一个矩阵化简为另一个矩阵时，第二个矩阵的每一行
都是第一个矩阵的各行的线性组合。
%</co:RowsOfEqMatsLinCombos>
\end{corollary}

% The proof 
% uses induction. % \appendrefs{mathematical induction}\spacefactor=1000 %
% Before we proceed, here is an outline of the argument.
% For the base step, we
% will verify that the proposition is true when reduction 
% can be done in zero row operations.
% For the inductive step, we will 
% argue that if being able to reduce the first matrix to the second in some
% number $t\geq 0$ of operations implies that each row of the second is a linear
% combination of the rows of the first, then being able to reduce the first to
% the second in $t+1$ operations implies the same thing.
% Together these prove the result because  
% the base step shows that it is true in the zero operations case,
% and then the inductive step
% implies that it is true in the one operation case, and then the inductive step
% applied again gives that it is therefore true for two operations, etc.

\begin{proof}
%<*pf:RowsOfEqMatsLinCombos0>
对于任何两个可以互相化简的矩阵 $A$ 与~$B$，存在某个
把其中一个变为另一个的最少行变换次数。
我们对该次数进行归纳。

在基础步骤中，即可以用零次化简操作从一个矩阵变到另一个矩阵时，
这两个矩阵相等。
于是 $B$ 的每一行显然是 $A$ 的各行的组合：
$\vec{\beta}_i
  =0\cdot\vec{\alpha}_1+\cdots+1\cdot\vec{\alpha}_i+\cdots+0\cdot\vec{\alpha}_m$。
%</pf:RowsOfEqMatsLinCombos0>

%<*pf:RowsOfEqMatsLinCombos1>
归纳步骤假设归纳假设：~即当 $k\geq 0$ 时，
任何能由 \( A \) 经 \( k \) 次或更少操作得到的矩阵，
其各行都是 $A$ 的各行的线性组合。
考虑矩阵~$B$，使得把 \( A \) 化简为~\( B \)
需要$k+1$次操作。
在该化简中有一个倒数第二个矩阵~$G$，
即 $A\longrightarrow\cdots\longrightarrow G\longrightarrow B$。
归纳假设适用于这个 \( G \)，
因为它距 \( A \) 只有$k$步。
也就是说，\( G \) 的每一行
都是 \( A \) 的各行的线性组合。
%</pf:RowsOfEqMatsLinCombos1>

%<*pf:RowsOfEqMatsLinCombos2>
我们将验证~$B$ 的各行是~$G$ 的
各行的线性组合。
于是，线性组合引理（\nearbylemma{lm:LinearCombinationLemma}）
适用，它表明~$B$ 的各行是~$A$ 的
各行的线性组合。

如果把 \( G \) 变为~\( B \) 的行变换是对换，那么
$B$ 的行不过是 $G$ 的行的重新排列，$B$ 的每一行
都是 $G$ 的各行的线性组合。
如果把 $G$ 变为~$B$ 的操作是用标量~$c\rho_i$
乘某一行，
那么 $\vec{\beta}_i=c\vec{\gamma}_i$，其余各行不变。
最后，如果行变换是把一行的倍数加到
另一行上 $r\rho_i+\rho_j$，
那么 $B$ 中只有第~$j$行与~$G$ 的对应行不同，且
$\vec{\beta}_j=r\gamma_i+\gamma_j$，
它确实是 $G$ 的各行的线性组合。
%</pf:RowsOfEqMatsLinCombos2>

由于我们既证明了基础步骤又证明了归纳步骤，
根据数学归纳法原理，该命题得证。
\end{proof}

现在我们认识到，高斯消元法构造的是
行的线性组合。
但它的目标当然是终止于阶梯形，因为那是
线性方程组的一种特别基本的形态，
它把变量分离了出来。
例如，在这个矩阵
\begin{equation*}
  R=\begin{mat}[r]
    2  &3  &7  &8  &0  &0  \\
    0  &0  &1  &5  &1  &1  \\
    0  &0  &0  &3  &3  &0  \\
    0  &0  &0  &0  &2  &1
  \end{mat}
\end{equation*}
中，$x_1$ 已从 $x_5$ 的方程中被消去。
也就是说，高斯消元法使 $x_5$ 的行在某种意义上
独立于 $x_1$ 的行。

下面的结果使这个直觉变得精确。
我们有时把高斯消元法称作高斯消去（Gaussian elimination）。
它所消去的是各行之间的线性关系。

\begin{lemma}      \label{le:EchFormNoLinCombo}
%<*lm:EchFormNoLinCombo>
在阶梯形矩阵中，
没有哪个非零行是其余非零行的线性组合。
%</lm:EchFormNoLinCombo>
\end{lemma}

\begin{proof}
%<*pf:EchFormNoLinCombo0>
设 $R$ 是一个阶梯形矩阵，考察它的非 \( \vec{0} \) 行。
首先注意到，
如果把一行写成
其余各行的组合
$\vec{\rho}_i=c_1\vec{\rho}_1+\cdots+c_{i-1}\vec{\rho}_{i-1}+
               c_{i+1}\vec{\rho}_{i+1}+\cdots+c_m\vec{\rho}_m$
那么可以把该方程改写为
\begin{equation*}
   \vec{0}=c_1\vec{\rho}_1+\cdots+c_{i-1}\vec{\rho}_{i-1}+c_i\vec{\rho}_i+
               c_{i+1}\vec{\rho}_{i+1}+\cdots+c_m\vec{\rho}_m
  \tag{$*$}
\end{equation*}
其中并非所有系数都为零；具体地说，$c_i=-1$。
反过来也成立：~给定方程~($*$)，若某个 $c_i\neq 0$，则通过把
$c_i\vec{\rho}_i$ 移到左边并除以 $-c_i$，
我们可以把 $\vec{\rho}_i$ 表示成其余各行的组合。
因此，如果我们证明在~($*$) 中所有系数都是~$0$，
就证明了该定理。
为此我们对行号~$i$ 使用归纳法。
%</pf:EchFormNoLinCombo0>

%<*pf:EchFormNoLinCombo1>
基础情形是第一行~$i=1$
（如果没有这样的非零行，从而 $R$ 是零矩阵，则
引理平凡地成立）。
设 $\ell_i$ 是第~$i$行首主元所在的列号。
考察每一行位于第~$\ell_1$列的元素。
方程~($*$) 给出
\begin{equation*}
  0=c_1r_{1,\ell_1}+c_2r_{2,\ell_1}+\cdots+c_mr_{m,\ell_1}
  \tag{$**$}
\end{equation*}
该矩阵是阶梯形的，所以
第一行之后的每一行在该列的元素都是零：
$r_{2,\ell_1}=\cdots=r_{m,\ell_1}=0$。
于是方程~($**$) 表明 $c_1=0$，因为 $r_{1,\ell_1}\neq 0$，
它是该行的首主元。
%</pf:EchFormNoLinCombo1>

%<*pf:EchFormNoLinCombo2>
归纳步骤与基础步骤基本相同。
再次考察方程~($*$)。
我们将证明：如果对每个行标 $i\in\set{1,\ldots,k}$
系数 $c_i$ 都是 $0$，
那么 $c_{k+1}$ 也是 $0$。
我们关注来自第~$\ell_{k+1}$列的元素。
\begin{equation*}
  0=c_1r_{1,\ell_{k+1}}+\cdots+c_{k+1}r_{k+1,\ell_{k+1}}+\cdots+c_mr_{m,\ell_{k+1}}
\end{equation*}
由归纳假设，$c_1$、\ldots、$c_k$ 全是
$0$，于是上式化归为方程
$0=c_{k+1}r_{k+1,\ell_{k+1}}+\cdots+c_mr_{m,\ell_{k+1}}$。
该矩阵是阶梯形的，所以元素
$r_{k+2,\ell_{k+1}}$、\ldots、$r_{m,\ell_{k+1}}$ 全是~$0$。
于是 $c_{k+1}=0$，因为 $r_{k+1,\ell_{k+1}}\neq 0$，它是首主元。
%</pf:EchFormNoLinCombo2>
\end{proof}

有了这些，我们就可以证明高斯–若尔当消元的
最终产物是唯一的。

\begin{theorem}
\label{th:ReducedEchelonFormIsUnique}
%<*th:ReducedEchelonFormIsUnique>
每个矩阵都行等价于唯一的简化阶梯形矩阵。
%</th:ReducedEchelonFormIsUnique>
\end{theorem}

\begin{proof} \cite{Yuster}
%<*pf:ReducedEchelonFormIsUnique0>
固定行数 \( m \)。
我们将对列数 \( n \) 进行归纳。

基础情形是矩阵只有 \( n=1 \) 列。
如果它是零矩阵，那么它的阶梯形是零矩阵。
如果它含有非零元素，那么当该矩阵被化为
简化阶梯形时，它必含至少一个非零元素，而这个元素
必是第一行中的一个 \( 1 \)。
无论哪种情形，它的简化阶梯形都是唯一的。
%</pf:ReducedEchelonFormIsUnique0>

%<*pf:ReducedEchelonFormIsUnique1>
归纳步骤假设 \( n>1 \)，并且所有列数少于~\( n \) 的
\( m \)~行矩阵都有唯一的简化阶梯形。
考虑一个 \( \nbym{m}{n} \) 矩阵 \( A \)，并设
\( B \) 与 \( C \) 是由 \( A \) 得到的两个简化阶梯形矩阵。
我们将证明这两个矩阵必相等。
%</pf:ReducedEchelonFormIsUnique1>

%<*pf:ReducedEchelonFormIsUnique2>
设 \( \hat{A} \) 是由 \( A \) 的前 \( n-1 \) 列
组成的矩阵。
注意到
% if an $n$~column matrix is in reduced echelon form then any initial
% set of its columns is also in reduced echelon form, so \( \hat{A} \)
% is in reduced echelon form. 
任何把 \( A \) 化为简化
阶梯形的行变换序列也会把 \( \hat{A} \) 化为简化阶梯形。
由归纳假设，\( \hat{A} \) 的这个简化阶梯形
是唯一的，所以如果 \( B \) 与 \( C \) 不同，那么差异必
出现在第~\( n \) 列。
%</pf:ReducedEchelonFormIsUnique2>

我们通过证明这两者不可能只在该列不同，
来完成归纳步骤和整个论证。
%<*pf:ReducedEchelonFormIsUnique3>
考虑以 \( A \) 为系数矩阵的一个齐次方程组。
\begin{equation*}
  \begin{linsys}{4}
    a_{1,1}x_1  &+  &a_{1,2}x_2  &+  &\cdots  &+  &a_{1,n}x_n  &=  &0  \\
    a_{2,1}x_1  &+  &a_{2,2}x_2  &+  &\cdots  &+  &a_{2,n}x_n  &=  &0  \\
              &&&&&&&\vdotswithin{=}  \\
    a_{m,1}x_1  &+  &a_{m,2}x_2  &+  &\cdots  &+  &a_{m,n}x_n  &=  &0  
  \end{linsys}
  \tag{$*$}
\end{equation*}
根据定理~One.I.\ref{th:GaussMethod}
%the first theorem of this chapter 
该方程组的解集与 $B$ 的方程组
\begin{equation*}
  \begin{linsys}{4}
    b_{1,1}x_1  &+  &b_{1,2}x_2  &+  &\cdots  &+  &b_{1,n}x_n  &=  &0  \\
    b_{2,1}x_1  &+  &b_{2,2}x_2  &+  &\cdots  &+  &b_{2,n}x_n  &=  &0  \\
               &&&&&&&\vdotswithin{=}  \\
    b_{m,1}x_1  &+  &b_{m,2}x_2  &+  &\cdots  &+  &b_{m,n}x_n  &=  &0  
  \end{linsys}
  \tag{$**$}
\end{equation*}
的解集相同，也与 $C$ 的
\begin{equation*}
  \quad
  \begin{linsys}{4}
    c_{1,1}x_1  &+  &c_{1,2}x_2  &+  &\cdots  &+  &c_{1,n}x_n  &=  &0  \\
    c_{2,1}x_1  &+  &c_{2,2}x_2  &+  &\cdots  &+  &c_{2,n}x_n  &=  &0  \\
               &&&&&&&\vdotswithin{=}  \\
    c_{m,1}x_1  &+  &c_{m,2}x_2  &+  &\cdots  &+  &c_{m,n}x_n  &=  &0  
  \end{linsys}
  \tag{$\mathord{*}\mathord{*}\mathord{*}$}
\end{equation*}
%</pf:ReducedEchelonFormIsUnique3>
%<*pf:ReducedEchelonFormIsUnique4>
设 \( B \) 与 \( C \) 只在第~\( n \) 列不同，并假设
它们在第~\( i \) 行不同。
用 ($**$) 的第~\( i \) 行减去 ($\mathord{*}\mathord{*}\mathord{*}$) 的
第~\( i \) 行，
得到方程
\( (b_{i,n}-c_{i,n})\cdot x_n=0 \)。
我们已假设 \( b_{i,n}\neq c_{i,n} \)，于是得到 $x_n=0$。
因此 $x_n$ 不是自由变量，
于是在 ($**$) 和~($\mathord{*}\mathord{*}\mathord{*}$) 中
第 \( n \) 列含有某行的首主元，
因为在阶梯形矩阵中任何
不含首主元的列都对应于
一个自由变量。

但现在，由于 \( B \) 与 \( C \) 在
前 \( n-1 \)~列相等，根据
简化阶梯形的定义，它们在
第 \( n \) 列的首主元位于同一行。
而且，这两个首主元都必须是 $1$，
并且必须是该列中仅有的非零元素。
因此 \( B=C \)。
%</pf:ReducedEchelonFormIsUnique4>
\end{proof}

我们曾问过：
一个线性方程组的任意两个阶梯形版本
是否有相同数目的自由变量，如果有，
它们是否恰好是同样的变量？
有了前面的结果，我们可以对这两个问题都回答“是。”
不存在这样的线性方程组，使得
比如我们用一种方式应用高斯消元法得到 $y$ 与 $z$
是自由变量，而用另一种方式应用却得到 $y$ 与 $w$ 是自由变量。

在证明之前，
回顾一下自由变量与参数的区别。
这个方程组
\begin{equation*}
  \begin{linsys}{3}
    x &+ &y &  &  &= &1 \\ 
      &  &y &+ &z &= &2  
  \end{linsys}
\end{equation*}
有一个自由变量~$z$，因为它是唯一不引领一行的变量。
我们习惯用自由变量来参数化：$y=2-z$，$x=-1+z$，
但我们也可以用另一个变量参数化，例如
$z=2-y$，$x=1-y$。
所以参数的集合并不唯一，唯一的
是自由变量的集合。

\begin{corollary}
如果从一个初始线性方程组出发，我们用高斯消元法得到两个
不同的阶梯形方程组，那么这两个方程组有相同的自由变量。
\end{corollary}

\begin{proof}
前面的结果说简化阶梯形是唯一的。
从任一阶梯形版本出发，
向上消元即可得到简化阶梯形，
所以方程组的任一阶梯形版本与简化阶梯形版本
有相同的自由变量。
\end{proof}

我们以一段回顾来收尾。
在高斯消元法中，我们从一个矩阵出发，
然后导出一系列其他的矩阵。
如果我们能从一个矩阵导出另一个，我们就称这两个矩阵是相关的。
这种关系是一种等价关系，%\appendrefs{equivalence relation} 
称为行等价，
于是它把全体矩阵的集合划分为一些行等价类。
\begin{center}
  \includegraphics{gr/mp/ch1.30}
\end{center}
（图中所示的类中有无穷多个矩阵，但限于篇幅
我们只展示两个。）
我们已经证明，每个行等价类中
有且仅有一个简化阶梯形矩阵。
%<*ReducedEchelonFormIsCanonicalForm>
所以简化阶梯形是
行等价的一种规范形\appendrefs{canonical representatives}\spacefactor=1000
\index{canonical form!for row equivalence}\index{representative!for row equivalence classes}
：
简化阶梯形矩阵正是
这些类的代表。
%</ReducedEchelonFormIsCanonicalForm>
\begin{center}
  \includegraphics{gr/mp/ch1.31}
\end{center}

这里的思想是：理解一个数学情境的一种方式，
是能够对可能发生的各种情形加以分类。
这是本书的一个主题，
我们已经见过它好几次了。
我们曾把线性方程组的解集分为无元素、一个元素
与无穷多元素三种情形。
我们也曾把方程个数与未知量个数相同的线性方程组
分为非奇异与奇异两种情形。

这里，在我们考察行等价时，我们知道全体矩阵的
集合分裂成一些行等价类，
而且我们现在有办法确切指出每一个这样的类\Dash
我们可以把一个类中的矩阵看成是由该类中
唯一的简化阶梯形矩阵经行变换导出的。

用更操作性的说法，简化阶梯形的唯一性
使我们能够通过把关于类的问题
转化为关于代表的问题
来回答它们。
例如，正如本节开头所允诺的，现在我们可以
判定一个矩阵能否由另一个矩阵经行化简导出。
我们对两者都应用高斯–若尔当程序，看它们
是否给出相同的简化阶梯形。

\begin{example}  \label{ex:MatsNotRowEq}
这两个矩阵不是行等价的
\begin{equation*}
  \begin{mat}[r]
    1  &-3  \\
   -2  &6
  \end{mat}
  \qquad
  \begin{mat}[r]
    1  &-3  \\
   -2  &5
  \end{mat}
\end{equation*}
因为它们的简化阶梯形不相等。
\begin{equation*}
  \begin{mat}[r]
    1  &-3  \\
    0  &0
  \end{mat}
  \qquad
  \begin{mat}[r]
    1  &0   \\
    0  &1
  \end{mat}
\end{equation*}
\end{example}

\begin{example}
任何非奇异的 \( \nbyn{3} \) 矩阵经高斯–若尔当消元都化为
下面这个矩阵。
\begin{equation*}
    \begin{mat}[r]
      1  &0  &0 \\
      0  &1  &0 \\
      0  &0  &1
    \end{mat}
\end{equation*}
\end{example}

\begin{example} \label{ex:RowEqClassTwoTwoMats}
我们可以通过列出所有可能的
简化阶梯形矩阵来描述全部的类。
任何 $\nbyn{2}$ 矩阵都属于下面这些之一：~与这个矩阵
行等价的矩阵类，
\begin{equation*}
  \begin{mat}[r]
     0  &0  \\
     0  &0
  \end{mat}
\end{equation*}
与这种类型的某个矩阵行等价的矩阵构成的
无穷多个类
\begin{equation*}
  \begin{mat}
     1  &a  \\
     0  &0
  \end{mat}
\end{equation*}
其中 \( a\in\Re \)（包括 $a=0$），
与这个矩阵行等价的矩阵类，
\begin{equation*}
  \begin{mat}[r]
     0  &1  \\
     0  &0
  \end{mat}
\end{equation*}
以及与这个矩阵行等价的矩阵类
\begin{equation*}
  \begin{mat}[r]
     1  &0  \\
     0  &1
  \end{mat}
\end{equation*}
（这就是非奇异的 $\nbyn{2}$ 矩阵构成的类）。
\end{example}



